% Part: counterfactuals
% Chapter: minimal-change-semantics
% Section: introduction

\documentclass[../../../include/open-logic-section]{subfiles}

\begin{document}

\olfileid{cnt}{min}{int}

\olsection{جان پہچان}

سٹالنیکر تے لوئس نے ’جے تیلی رگڑی جاندی، تاں بل پَیندی‘ ورگے خلافِ واقعہ
شرطیاں دیاں تشریحاں پیش کیتیاں۔ ایہ تجویزاں ایہہ سمجھاؤن لئی سن کہ ایسے جملیاں
دے سچ ہون دیاں شرطاں کیویں سمجھیاں جان۔ دوناں دا بنیادی خیال ایہ اے۔ خلافِ
واقعہ شرطیے دا سچ پرکھن لئی اوہ ممکن جہان ویکھنے پین گے جیہڑے اصل جہان توں
اینے کو وکھرے ہون کہ مقدم سچ ہو جاوے، تے ایہ فرق گھٹ توں گھٹ ہووے۔
جے تالی اوہناں ممکن جہاناں وچ سچ اے، تاں خلافِ واقعہ شرطیہ سچ اے۔
مثلاً منّ لو میرے ہتھ وچ اک تیلی تے تیلیاں دا اک پیکٹ اے۔ اصل جہان وچ
میں بس اوہناں نوں ویکھدا آں تے سوچدا آں کہ جے میں تیلی رگڑاں تاں کی ہووے۔
اصل جہان توں گھٹ توں گھٹ تبدیلی والا جہان جتھے میں تیلی رگڑدا آں، اوہ اے
جتھے میں کرن دا فیصلہ کردا آں تے تیلی رگڑ دیندا آں۔ ایہ تبدیلی گھٹ توں گھٹ
ایس لئی اے کہ ہور کجھ نہیں بدلدا۔ میں نال ہوا وچ چھلانگ نہیں لاندا، تیلی
رگڑن نال میرے والاں نوں اگ نہیں لگدی، میریاں اُنگلاں دا زور اچانک نہیں
مکدا، تے نال ای کسے لُک کے سپر سوکر نال میرے اُتے پانی نہیں سُٹیا، وغیرہ۔
اوس ہور ممکن صورت وچ تیلی بل پَیندی اے۔ ایس لئی ایہ سچ اے کہ جے میں
تیلی رگڑدا، تاں اوہ بل پَیندی۔

ایس سوکھی سمجھ آؤن والی تشریح نوں خلافِ واقعہ شرطیاں دیاں منطقاں دی رسمی
معنیات نال جوڑ سکدے آں۔ لوئس نے خلافِ واقعہ شرطیے لئی نشان ’$\boxright$‘
پیش کیتا، تے سٹالنیکر نے نشان~’$>$‘ ورتیا۔ اسیں~$\cif$ ورتاں گے تے اینوں
قضیاتی منطق وچ دو-مقامی رابطے دے طور اُتے شامل کراں گے۔ سو !!{formula}s جیہڑے $!A \lif
!B$ دی شکل دے نیں، اوہناں توں علاوہ !!{formula}s جیہڑے~$!A \cif !B$ دی
شکل دے نیں، وی نیں۔ رسمی معنیات، امکانی منطق دی تعلقی معنیات وانگ، اوہناں نمونیاں
اُتے ٹکی اے جیہناں وچ !!{formula}s جہاناں اُتے پرکھے جاندے نیں۔ پورا اترن دی
شرط جیہڑی $\mSat{M}{!A \cif !B}[w]$ دی تعریف کردی اے، $\mSat{M}{!A}[w']$ تے $\mSat{M}{!B}[w']$ دے لحاظ نال،
کجھ (ہور) جہاناں~$w'$ لئی دتی جاندی اے۔ کیہڑے $w'$؟ سمجھ مطابق،~$w$ دے سبھ توں نیڑے اوہ جہان
جیہناں لئی~$\mSat{M}{!A}[w']$ اے۔ ایس لئی نمونے وچ ’نیڑتا‘ دا تعلق وی
شامل ہونا چاہیدا اے۔

لوئس نے خلافِ واقعہ حالتاں نوں نقشے نال وکھاؤن دا اک سمجھ آؤن والا
طریقہ پیش کیتا۔ ہر ممکن جہان اک دوجے اندر سمائے ہوئے گولیاں دے سیٹ دا مرکز
اے۔ گولیاں اندر ہور جہان نیں، تے اسیں گولیاں نوں سانجھے مرکز والے دائریاں
وانگ بنا کے وکھاندے آں۔ دو گولیاں دے وچلے جہان مرکز والے جہان توں سارے
اکّو جِنّے نیڑے نیں۔ اندرلے گولے وچلے جہان ہور نیڑے، تے آلے دوالے دے گولے
وچلے جہان ہور دُور نیں۔
\begin{center}
\begin{tikzpicture}[scale=.7]\small
  \spheresystem{5}
  \draw (0,0) node {$w$};
  \propositionintersect{0}{4}{40}{3.5}
  {\tikzset{proposition/.append={smooth,tension=1.4}}
  \proposition[shift={(1.6,-1)}]{90}{1}{30}{4}}
  \path (1.5,3) node[below] {$\formula{B}$};
  \path (3,0) node {$\formula{A}$};
\end{tikzpicture}
\end{center}
سبھ توں نیڑلے $!A$ والے جہان اوہ جہان $w'$ نیں جتھے~$!A$ پورا اتردا اے
تے جیہڑے مرکز والے جہان~$w$ دے آلے دوالے سبھ توں چھوٹے گولے وچ نیں
(سُرمئی حصہ)۔ سمجھ مطابق، $!A \cif !B$،~$w$ اُتے پورا اتردا اے جے $!B$
سارے سبھ توں نیڑلے $!A$ والے جہاناں اُتے سچ ہووے۔

\end{document}
